\documentclass[12pt]{article}
\usepackage{amsmath}
\usepackage{amssymb}
\usepackage[margin=1in]{geometry}

\title{Can a simple model explain economics?}
\author{Ludovico Furlanetto}
\date{July 2026}

\begin{document}
\maketitle

\begin{abstract}
Firm growth shows two robust anomalies: the volatility of a firm's growth rate falls
with size far more slowly than independence would predict, $\sigma(S)\sim S^{-\beta}$
with $\beta\approx0.15$--$0.20$ rather than $1/2$, and the growth-rate distribution is
sharply peaked and fat-tailed rather than Gaussian. The standard explanations place the
source inside each firm, in a heavy-tailed composition of internal sub-units, a granular
picture that Moran et al.\ test against firm data and find fails; the missing
ingredient, they argue, is correlations that grow with a firm's size. We ask instead
whether a dynamical model of interacting firms can generate both facts on its own, without
external noise.

We model the economy as a Generalized Lotka--Volterra system on a scale-free random graph
with the standard dilute couplings, in
which the interactions act on a firm's size relative to the average firm rather than on
its absolute size. This makes the dynamics
scale-invariant, so the economy grows exponentially, while in the strongly disordered
regime its composition never settles but keeps churning chaotically. In this persistently
fluctuating state both anomalies emerge together as stationary properties of the
interactions, from a single deterministic ingredient, the disordered couplings, with no
external noise. The size--variance relation is a clean power law
with the empirical sign, $\beta\approx0.45$ at our largest sizes. Its higher moments
reproduce the anomalous, order-dependent
multiscaling the granular picture cannot: normalized by the first,
$\zeta_q/\zeta_1\approx1,\,1.84,\,2.51,\,3.02$ lies on the empirical
$1,\,2.0,\,2.6,\,2.9$, converged in system size; the growth-rate distribution is
symmetric (Bowley $\approx0$) and fat-tailed; and the heavy tails do not merely persist
among the largest firms but strengthen with size, the excess kurtosis rising with size in
every realization where an aggregate of independent parts would approach a Gaussian. Two
ablations locate the origins. Switching off the disorder removes the fluctuations
entirely; swapping the graph shows the size-conditioned statistics to be generated by its
degree heterogeneity --- under the dilute normalization a firm's noise variance grows with
its connectivity, on homogeneous graphs the multiscaling vanishes exactly, and only the
scale-free graph reproduces the empirical curve. The network, not the single firm, is
where the size-growing correlations live.
\end{abstract}

\end{document}
